% authority: science_draft
% claim_status: draft/control scoped obstruction
\documentclass[11pt]{article}
\usepackage[margin=1in]{geometry}
\usepackage{amsmath,amssymb,amsthm,mathtools}
\usepackage{booktabs,array,longtable,enumitem,microtype}
\usepackage[hidelinks]{hyperref}
\setlist{nosep}

\newcommand{\A}{\mathcal A}
\newcommand{\R}{\mathbb R}
\newcommand{\dd}{\mathrm d}
\newcommand{\Gam}{\Gamma_{\tau}}
\newcommand{\sigmaL}{\sigma_{L_D}}
\newcommand{\Char}{\operatorname{Char}}
\newtheorem{definition}{Definition}
\newtheorem{proposition}{Proposition}
\newtheorem{theorem}{Theorem}
\newtheorem{lemma}{Lemma}

\title{V22 P4--T02: Hyperbolicity, Universality, Robustness,\
and Hard-Fail Screen for the B1 Six-Transport Candidate}
\author{The \AE ther-Flow Research Project}
\date{2026-08-09}

\begin{document}
\maketitle

\begin{abstract}
This draft/control packet stress-tests the exact P4--T01 reduced source
principal structure.  The fixed diagonal system has a precise positive result:
with respect to the proposal-only source label $\tau$, it is symmetric
hyperbolic, its field-fiber symmetrizer is the identity, and
$P_D=\prod_A k(X_A)$ is hyperbolic on the common algebraic component
$\Gam=\bigcap_A\{k:k(X_A)>0\}$.  This is source PDE structure, not physical
causality.  The candidate nevertheless fails the P4--T02 continuation
predicate.  Its reduced characteristic set is the union of six distinct
hyperplanes; no nonzero quadratic polynomial, hence no nondegenerate
Lorentzian quadratic, vanishes on that union.  No physical selector or quotient
identifies the branches, and the declared finite P7 sectors have no continuum
map into the B1 equations.  Moreover, an arbitrarily small, real, source-local
cross-channel mutation preserving the selected background gives
discriminant $-16\epsilon^2$ at a branch crossing and therefore complex
characteristic speeds for every $\epsilon\ne0$.  The fixed law remains
hyperbolic, but the broader B1 family lacks a source-derived rule protecting
that diagonal form.  We therefore trigger the preregistered HF04 and
family-level HF05 conditions, terminate B1 append-only, keep the B2 fallback
inactive, and require a fresh selector packet.  This is a candidate-scoped
obstruction, not a global no-go theorem, Gate B verdict, effective metric, or
future-extension closure.
\end{abstract}

\section{Inherited candidate and exact scope}

Let $\A$ be the proposal-only four-dimensional source arena, let
$D=(\tau,X_1,\ldots,X_6)$ be the fixed P3--T02 coefficient datum, and let
$q=(q^1,\ldots,q^6)$.  P4--T01 lawfully reduced the fixed-D tangent equations
to
\begin{equation}
 L_Du=(X_1u^1,\ldots,X_6u^6),
 \qquad
 \sigmaL(k)=\operatorname{diag}(k(X_1),\ldots,k(X_6)),
 \label{eq:fixed}
\end{equation}
with zero algebraic constraints, differential constraints, internal field
gauge, and equation redundancies.  In the declared source chart,
\begin{align}
 X_1&=\partial_0+\partial_1,&
 X_2&=\partial_0+\partial_2,&
 X_3&=\partial_0+\partial_3,\nonumber\\
 X_4&=\partial_0+\partial_1+\partial_2,&
 X_5&=\partial_0+\partial_2+2\partial_3,&
 X_6&=\partial_0+2\partial_1+\partial_3.
 \label{eq:vectors}
\end{align}
Every channel obeys $\dd\tau(X_A)=1$.  The reduced polynomial is
\begin{align}
 P_D(k)&=\prod_{A=1}^{6}k(X_A)\nonumber\\
 &= (\omega+\kappa_1)(\omega+\kappa_2)(\omega+\kappa_3)
 (\omega+\kappa_1+\kappa_2)
 (\omega+\kappa_2+2\kappa_3)
 (\omega+2\kappa_1+\kappa_3).
 \label{eq:poly}
\end{align}
The six factors are pairwise nonassociate and globally simple.  Directional
crossings are real and retained; every pair crosses and the maximum nonzero
symbol corank is four.

This packet asks two different questions and does not conflate them:
\begin{enumerate}[leftmargin=*]
  \item Is the exact fixed source evolution hyperbolic and well posed relative
        to its declared source label?
  \item Does that result supply one stable, universal, physically typed
        Lorentzian causal structure with adequate family-level robustness?
\end{enumerate}
The first answer is positive.  The second is negative in the exact scope below.

\section{Exact fixed-law source hyperbolicity}

Write $X_A=\partial_\tau+v_A^j\partial_j$.  Then the fixed evolution is
\begin{equation}
 \partial_\tau u+\sum_{j=1}^{3}A^j\partial_j u=0,
 \qquad
 A^j=\operatorname{diag}(v_1^j,\ldots,v_6^j).
 \label{eq:evolution}
\end{equation}

\begin{theorem}[Fixed source symmetric hyperbolicity]
Equation \eqref{eq:evolution} is symmetric hyperbolic relative to the
proposal-only source label $\tau$.  The time matrix is $A^0=I_6$, every
$A^j$ is real diagonal and symmetric, and $H=I_6$ is a positive field-fiber
symmetrizer.  The spatial symbol has the fixed canonical eigenbasis for every
$\kappa$, so its eigenvector condition number is one, including at repeated
directional roots.
\end{theorem}

\begin{proof}
The matrices in \eqref{eq:evolution} are diagonal by inspection of
\eqref{eq:fixed}.  Thus $H A^0=I_6$ is positive definite and each $H A^j$
is symmetric.  The standard channel basis diagonalizes every spatial symbol
$A(\kappa)=\sum_j A^j\kappa_j$ simultaneously.  Repeated eigenvalues therefore
do not destroy diagonalizability or condition the eigenbasis.
\end{proof}

The identity $H$ acts on six field amplitudes.  It is not a bilinear form on
$T\A$ or $T^*\A$, does not contract source derivatives, and does not define a
source-spacetime metric.  Under a field-frame change it transforms as a
field-energy form.  Treating it as $g_{\rm eff}$ would be a type error.

\begin{proposition}[Polynomial hyperbolicity and common source component]
The polynomial \eqref{eq:poly} is hyperbolic with respect to
$h=\dd\tau$.  The component containing $h$ is
\begin{equation}
 \Gam=\bigcap_{A=1}^{6}\{k\in T_x^*\A:k(X_A)>0\}.
 \label{eq:gamma}
\end{equation}
It is an open convex polyhedral source-cotangent cone.
\end{proposition}

\begin{proof}
Every factor satisfies $h(X_A)=1$.  For arbitrary $k$, the roots of
\begin{equation}
 \lambda\longmapsto P_D(k+\lambda h)
 =\prod_A\bigl(k(X_A)+\lambda h(X_A)\bigr)
\end{equation}
are $\lambda_A=-k(X_A)$ and are real.  The sign component containing $h$ is
exactly the intersection of the six positive factor half-spaces, which proves
\eqref{eq:gamma} and convexity.
\end{proof}

This establishes source evolution orientation and algebraic PDE hyperbolicity.
It does not establish a physical clock, light cone, Lorentzian signature, or
time orientation.  Those require a lawful source-to-world interpretation that
the B1 packet does not supply.

\section{No common Lorentzian characteristic quadric}

Let $\ell_A(k)=k(X_A)$.  The reduced characteristic variety is
\begin{equation}
 \Char(L_D)_x=\bigcup_{A=1}^{6}\{[k]:\ell_A(k)=0\}.
 \label{eq:char}
\end{equation}

\begin{theorem}[No common characteristic quadric]
There is no nonzero homogeneous quadratic polynomial $Q$ on $T_x^*\A$ that
vanishes on the full union in \eqref{eq:char}.  In particular, the reduced
characteristic set is not the null set of one nondegenerate Lorentzian
quadratic form.
\end{theorem}

\begin{proof}
If a polynomial vanishes on the hyperplane $\ell_A=0$, then the linear
polynomial $\ell_A$ divides it.  Therefore a polynomial vanishing on all six
hyperplanes is divisible by $\prod_A\ell_A$.  P4--T01 proved the six factors
pairwise nonassociate, so this divisor has degree six.  A nonzero polynomial
of degree two cannot contain it.
\end{proof}

An independent exact calculation represents a general quadratic by its ten
monomial coefficients.  Evaluating it at 156 integer points on the six
hyperplanes gives a rational constraint matrix of rank ten.  Hence only the
zero quadratic vanishes on all six branches.  This is an independent finite
certificate of the same obstruction, not external replication.

The pair intersections also make the reduced characteristic variety singular:
if two distinct factors vanish, every term in $\dd P_D$ contains at least one
vanishing factor.  A source hyperbolicity component still exists, but the
branch-complete characteristic geometry is reducible and crossing-rich.  No
candidate-supplied quotient identifies these branches, because P4--T01 found
zero internal gauge and zero algebraic redundancy.  Discarding a branch would
change the reduced system rather than reconstruct a physical equivalence.

\section{Universality and declared-sector audit}

The six B1 channels have source velocities
\begin{equation}
 (1,0,0),\ (0,1,0),\ (0,0,1),\ (1,1,0),\ (0,1,2),\ (2,0,1).
 \label{eq:velocities}
\end{equation}
They share the noncharacteristic source covector $\dd\tau$ and the common
component \eqref{eq:gamma}, but they do not share one characteristic
hypersurface.  For $\kappa=(1,0,0)$, for example, the six $\omega$ roots are
\begin{equation}
 (-1,0,0,-1,0,-2),
\end{equation}
and for $\kappa=(1,2,3)$ they are
$(-1,-2,-3,-3,-8,-5)$.  A common Cauchy surface is therefore not a common
propagation characteristic.

The declared Position-A P7 source input supplies finite charge labels $Q$,
defect classes $D$, representation grades $M$, and invariant pairs $(Q,D)$.
It supplies no continuum sector equation, quotient tangent, or map assigning a
P7 sector to a B1 channel.  Conversely, the B1 component labels are source
scalars and are not licensed as physical matter sectors.  The complete
comparison is therefore:
\begin{center}
\small
\begin{tabular}{p{0.25\linewidth}p{0.29\linewidth}p{0.36\linewidth}}
\toprule
Object & Principal information & P4--T02 status\\
\midrule
B1 channels $q^1,\ldots,q^6$ & six distinct source factors & common source time, no physical sector typing\\
P7 charge label $Q$ & no continuum equation & fail closed: not mapped\\
P7 defect class $D$ & no continuum equation & fail closed: not mapped\\
P7 grade $M$ & no continuum equation & fail closed: not mapped\\
P7 invariant pair $(Q,D)$ & finite transition invariant only & fail closed: no common-principal proof\\
\bottomrule
\end{tabular}
\end{center}
Thus the candidate has neither universal declared-matter propagation nor a
physically derived generalized causal equivalence.  Calling the polyhedral
source component a universal physical cone would identify source labels with
physical causal semantics by fiat.

\section{Finite-variation countermodel}

The exact diagonal system is well conditioned, but P3--T04 preserved a
cross-law selection debt: the current ontology does not derive why the
diagonal coefficient pattern is protected among the general source-local
first-jet B1 coefficient class.  The following mutation tests that debt.

Let $Y=\partial_1+\partial_2$ and let $\bar q$ be the selected P3--T04
background.  For real $\epsilon$, change only the first two equations to
\begin{align}
 E_1^\epsilon(q)&=X_1q^1+\epsilon Y(q^2-\bar q^2),\nonumber\\
 E_2^\epsilon(q)&=X_2q^2-\epsilon Y(q^1-\bar q^1).
 \label{eq:mutation}
\end{align}
The added terms vanish identically at $q=\bar q$, so the background is
preserved.  The mutation is real, local, first order, and source-only.  It
uses no metric, target atlas, physical norm, or GR-fit coefficient.  Since $Y$
is spatial relative to $\tau$, the time matrix remains $I_6$.

\begin{theorem}[Arbitrarily small family-level loss of hyperbolicity]
At the branch-1/branch-2 crossing $\kappa=(1,1,0)$, the first two spatial
symbol entries of \eqref{eq:mutation} form
\begin{equation}
 B_\epsilon=
 \begin{pmatrix}1&2\epsilon\\-2\epsilon&1\end{pmatrix}.
 \label{eq:block}
\end{equation}
Its characteristic discriminant is $-16\epsilon^2$.  Thus every nonzero
$\epsilon$, however small, gives a complex-conjugate characteristic-speed
pair.
\end{theorem}

\begin{proof}
The characteristic polynomial is
$(\lambda-1)^2+4\epsilon^2$, whose roots are
$1\pm2i\epsilon$.  Equivalently, the quadratic discriminant is
$-16\epsilon^2<0$ for $\epsilon\ne0$.
\end{proof}

This theorem does not refute the immutable exact diagonal law: it changes that
law.  It refutes family-level finite-variation robustness under the
prospectively admitted B1 source-local coefficient class unless a new
source-derived restriction protects simultaneous diagonal coupling.  Such a
restriction is not present or adopted.

The exact sweep at
$\epsilon=10^{-5},10^{-4},10^{-3},10^{-2}$ reproduces a negative discriminant
at every nonzero value.  By contrast, diagonal coefficient variations that
preserve $\dd\tau(X_A)>0$ retain source symmetric hyperbolicity.  The failure
is therefore accurately localized to unprotected cross-channel variation at
branch crossings rather than misreported as fixed-law nonhyperbolicity.

\section{Refinement and conditioning}

P3--T03 proves a controlled upwind refinement limit for the exact fixed law.
For mode $(1,1,1)$, the maximum discrete factor errors at
$N=8,16,32,64$ are approximately
\begin{equation}
 7.2762,\quad3.6853,\quad1.8486,\quad0.9250,
\end{equation}
and decrease monotonically.  The limit preserves all six factors.  Hence
HF06 is not triggered for the exact fixed law.  The same result also proves
that refinement does not merge the factors into one quadratic cone: it
converges to the branch structure that fails GB03 and GB04.

At the fixed-law level, the canonical eigenbasis has condition number one and
directional multiplicities are controlled.  At the family level, the
countermodel \eqref{eq:mutation} shows that arbitrary small cross-channel
variation can leave the real spectrum.  This distinguishes numerical
convergence of one sealed discretization from robustness of the scientific
candidate family.

\section{Hidden-target and hard-fail screen}

No input uses a target atlas, source-cotangent metric, tetrad, coframe,
Lorentzian constitutive tensor, target detector, GR target value, or Einstein
action.  The field symmetrizer $H$ and source label $\tau$ are kept in their
declared types.  Therefore HF01, HF02, HF03, HF07, and HF08 are not triggered;
absence of those failures is hygiene evidence only.

\begin{longtable}{>{\raggedright\arraybackslash}p{0.12\linewidth}
                  >{\raggedright\arraybackslash}p{0.20\linewidth}
                  >{\raggedright\arraybackslash}p{0.58\linewidth}}
\toprule
Criterion & Status & Exact reason\\
\midrule
GB01 & pass, hygiene only & complete source provenance and no target atlas\\
GB02 & pass, source PDE only & reduced $P_D$ hyperbolic and fixed system symmetric hyperbolic\\
GB03 & fail & no nonzero common quadratic characteristic polynomial and no physical selector\\
GB04 & fail & distinct B1 branches and no P7-to-B1 continuum sector map\\
GB05 & not established & no physical clock, rod, calibration, units, or scale\\
GB06 & partial source only & passive covariance and local sheaf gluing, not physical global gluing\\
GB07 & fail at family level & infinitesimal source-only cross-coupling creates complex speeds\\
GB08 & pass, hygiene only & no metric-equivalent primitive; field symmetrizer kept typed\\
\bottomrule
\end{longtable}

Two preregistered hard-fail triggers control the candidate disposition:
\begin{enumerate}[leftmargin=*]
  \item \texttt{HF04\_NONUNIVERSAL\_OR\_UNSELECTED\_MULTICONE}: the reduced
        source system has six distinct branches and no derived physical
        selector, quotient, or sector-common characteristic structure;
  \item \texttt{HF05\_INSTABILITY\_NONHYPERBOLICITY\_DEGENERACY}: the broader
        B1 coefficient family is not hyperbolicity-robust at crossings because
        \eqref{eq:mutation} gives complex speeds for arbitrarily small
        nonzero mutation.
\end{enumerate}
The phrase ``multicone'' here names the preregistered trigger.  The six B1
branches remain source characteristics; this packet does not promote them to
six physical cones.

\section{Disposition, freeze control, and next lawful route}

The result classification is \texttt{scoped\_obstruction}.  B1 has completed
the construct, continuum-control, linearize, reduce, and stress cycle.  It is
terminated and locally frozen under the exact tested assumptions.  The
termination record is append-only, the active-family slot is cleared, and
family slot 1 is consumed.

The preregistered B2 candidate
\texttt{CAND-V22-B2-P7-COMMON-PRINCIPAL-LIFT-V1} remains an inactive fallback.
It is not automatically activated.  Its missing continuum-lift descriptor,
sector equations, quotient data, common-principal target, no-target-import
receipt, and post-lift adequacy evaluation still fail closed.  P2--T03
therefore requires a fresh Theoretical Continuation Selector decision.  That
decision is a separate AgentJob and may select only one bounded B2 descriptor
or activation path while preserving B1's negative result.

P4--T03 is not unlocked because P4--T02 did not issue a qualifying continuation
verdict.  The three-family budget is not exhausted: B2 remains preregistered
and slot 3 remains reserved and unallocated.  The exact permitted conclusion
is that this B1 candidate fails under its declared source-family scope.  It is
not that all source extensions, the global ontology, or all \AE ther theories
fail.

\section{Distance-to-GR status}

\begin{longtable}{p{0.29\linewidth}p{0.62\linewidth}}
\toprule
Burden & Status after P4--T02 B1 screen\\
\midrule
Source ontology primitives & unchanged; B1 data remain proposal-only and B1 is locally frozen\\
Source equivalence $\mathrm{EqSrc}$ & passive source-chart covariance retained; no physical quotient added\\
$\mathrm{RetainH}$ and $\mathrm{GenH}$ & unchanged and not derived\\
$\mathrm{ObsLoc}_{\rm lc}$ and $\mathrm{Resp}_{\rm lc}$ & B1 token rank readiness inherited; no physical observer or universal sector map\\
$M_{\rm src}$ & not adopted or identified with physical spacetime\\
$g_{\rm eff}$ & blocked; no stable universal Lorentzian cone, scale, or metric exists in this packet\\
Matter coupling & blocked; finite P7 sectors have no B1 continuum principal map\\
Einstein equations & blocked\\
Finite-variation robustness & fixed-law source PDE positive; B1 family-level criterion fails by exact countermodel\\
Benchmark and Gate status & blocked; no protected verdict or external review requested\\
Current route freeze & B1 locally frozen; fresh B2 selector route remains open\\
Hard-fail status & HF04 and family-level HF05 trigger candidate termination\\
\bottomrule
\end{longtable}

The physical Distance-to-GR delta is zero.  A negative candidate disposition
sharpens route state but does not construct $g_{\rm eff}$ or promote a
downstream physical object.

\section{Executable and review evidence}

The primary standard-library implementation checks exact fixed-law roots,
source symmetric-hyperbolicity matrices, the hyperbolicity component,
quadratic divisibility obstruction, exact perturbation discriminants, and
refinement-factor convergence.  A separate implementation imports no primary
model.  It constructs a 156-by-10 rational quadratic-constraint matrix of
rank ten and independently reproduces the negative perturbation discriminants.

The Physicist-Mathematician child owns the PDE proofs, exact countermodel, and
quadratic obstruction.  The Physicist-Philosopher child owns the source versus
physical type boundary, sector non-identification, hard-fail scope, and global
no-go prohibition.  The parent reconciles the apparent tension between
fixed-law hyperbolicity and candidate failure: both are true because the
positive result is source-PDE-local while P4--T02 requires a stable universal
physical structure and family robustness.

Internal role separation and the independent implementation are not external
PDE review or independent scientific replication.  A validator or checkpoint
PASS is not proof authority, Gate authority, or physics promotion.

\section{Final statement}

The decisive P4--T02 B1 result is
\begin{center}
\fbox{\parbox{0.88\linewidth}{\centering
fixed source symmetric hyperbolicity\quad $+$\quad
candidate-scoped universal/robust Lorentzian obstruction.}}
\end{center}
B1 terminates under HF04 and family-level HF05.  B2 remains inactive pending a
fresh selector and its missing construction prerequisites.  No canonical
ontology edit, source-law adoption, physical time, physical cone, Lorentzian
signature, scale, effective metric, matter coupling, Einstein equation, Gate
B verdict, benchmark promotion, publication, push, global no-go theorem, or
completed derivation follows.

\end{document}
